\section{Horizontal poles over projective space}
\label{sec:horizontal-kernel-asymptotics}

When the base divisor is a hyperplane in projective space, the descended vanishing conditions of \cref{thm:horizontal-center-descent} turn the kernels into binomial tail probabilities. The large-deviation rate of these tails is the domination envelope of the original singularity upstairs, and the concentration of evaluation sections identifies the limiting measure for every smooth integration volume. In this setting the weak limit of \cref{cor:horizontal-center-kernel-limit} improves to convergence in total variation.

We specialize \cref{sec:horizontal-center-constraints} to $B=\mathbb P^q$ with homogeneous coordinates $[Z_0:\cdots:Z_q]$, $A=\mathcal O_B(k)$ and $Z=\{Z_0=0\}$, where $q,r,k$ are positive integers and $0<c<k$. Let $\omega_B$ be the Fubini--Study form on $B$, normalized by $\int_B\omega_B^q=1$, and give $A$ the Fubini--Study metric, whose curvature is $\theta=k\omega_B$. We keep the submersion $\pi:X\to B$ with connected $r$-dimensional fibers, and the selecting function $\varphi$ with the singularity \eqref{eq:horizontal-center-singularity} along the section $Y$ over $Z$. Write
\[
 R=\frac{|Z_0|^2}{\sum_{i=0}^q|Z_i|^2},\qquad
 a=\frac ck,\qquad D_c=\{R\geq a\},\qquad
 C_c=\frac{(k-c)^q}{q!}.
\]
Fix a smooth positive probability volume $\nu_X$ on $X$. Fiber integration gives $\nu_B=\pi_*\nu_X=w\,\omega_B^q$ with a smooth function $w>0$; put $M=\max_Bw$. Let $B_m^B$ be the kernel of the descended space $V_m$ for the norm on $B$ with zero continuous weight and volume $\nu_B$, and let $B_m^X$ be the kernel of the pulled-back space for the norm on $X$ with zero continuous weight and volume $\nu_X$. We shall see that $B_m^X=B_m^B\circ\pi$.

The selecting function $\varphi$ need not be $\pi^*\theta$-plurisubharmonic. Its domination envelope in this curvature background is
\[
 \mathcal P(\varphi)=
 \left(\sup\{u\in\PSH(X,\pi^*\theta):u\leq0,\ u\preceq\varphi\}\right)^*,
\]
where $u\preceq\varphi$ means that $u\leq\varphi+C$ for some constant $C$ (\cref{def:singularity-order}). The limit will be described by the function of $R$ defined by
\begin{equation}\label{eq:horizontal-projective-envelope}
 u_c(R)=
 \begin{cases}
 -k\left(a\log\dfrac aR+(1-a)\log\dfrac{1-a}{1-R}\right),&0<R<a,\\
 0,&a\leq R\leq1,\\
 -\infty,&R=0.
 \end{cases}
\end{equation}

\begin{theorem}\label{thm:horizontal-projective-kernels}\index{Bergman kernel!horizontal pole}
The envelope satisfies $\mathcal P(\varphi)=u_c\circ\pi$. With
$\ell_m=\lfloor mc\rfloor$ and $t_m=\max\{\ell_m-r,0\}$,
\[
 N_m=\dim V_m=\binom{km-t_m+q}{q}
       =C_cm^q+O(m^{q-1}).
\]
Thus the section spaces grow like $m^q$, and $q<\dim X$. For every $m\geq1$, in the extended-value sense,
\begin{equation}\label{eq:horizontal-kernel-lower}
 \frac1m\log B_m^X\geq\mathcal P(\varphi)-\frac{\log M}{m}.
\end{equation}
The functions $m^{-1}\log B_m^X$ converge to $\mathcal P(\varphi)$ in $L^1$ for every smooth positive volume on $X$, and locally uniformly on $X\setminus\pi^{-1}(Z)$. In the normalization of \cref{sec:partial-selected-space}, the kernel measures converge in total variation:
\begin{equation}\label{eq:horizontal-kernel-tv}
 \begin{split}
 \frac{q!}{m^q}B_m^B\nu_B&\longrightarrow\mu_c:=\mathbf1_{D_c}\theta^q,\\
 \frac{q!}{m^q}B_m^X\nu_X&\longrightarrow
       \Lambda_{\nu_X}:=\frac{k^q\mathbf1_{D_c}\circ\pi}{w\circ\pi}\nu_X.
 \end{split}
\end{equation}
Both limits have mass $(k-c)^q$, and $\pi_*\Lambda_{\nu_X}=\mu_c$. Dividing the unnormalized kernel measures by $N_m$ instead gives the probability limits $\mu_c/(k-c)^q$ and $\Lambda_{\nu_X}/(k-c)^q$, while normalization by $m^{q+r}$ gives zero in total variation.

Let $W_m$ be the space of \cref{thm:horizontal-center-descent}. For every sequence of intermediate spaces $W_m\subseteq U_m\subseteq V_m$, equipped with the restricted norms, the lower bound and all these convergences remain valid. The dimensions of these spaces are $C_cm^q+O(m^{q-1})$, and the dimension normalization then uses $\dim U_m$.
\end{theorem}

\subsection{The domination envelope}

We first compute the envelope on the base, using the affine chart $Z_0\ne0$. Set $z_i=Z_i/Z_0$, $s=|z|^2$, $l=k-c$, $s_0=l/c$ and
$C=k\log k-c\log c-l\log l$. In this chart the reference weight of $A$ is $\phi_A=k\log(1+s)$. Put
\[
 F_0(z)=
 \begin{cases}
 k\log(1+s),&s\leq s_0,\\
 l\log s+C,&s\geq s_0.
 \end{cases}
\]
As a function of $\log s$, $F_0$ is increasing and convex: its derivative increases to $l$ and then stays constant. Hence $F_0$ is plurisubharmonic. The two branches match to first order at $s=s_0$, and the tangent-line inequality gives $F_0\leq\phi_A$.
On a chart $Z_j\ne0$ with $j\geq1$, write $x=Z_0/Z_j$ and denote the remaining ratios by $\zeta$. Near $Z$ the corresponding weight is
\[
 F_j=c\log|x|^2+l\log(1+|\zeta|^2)+C.
\]
Therefore $F_0$ and the $F_j$ define a global plurisubharmonic metric weight $F$ on $A$, with a divisor pole of coefficient $c$ along $Z$ and a smooth remainder.

Let $v\leq\phi_A$ be a competing plurisubharmonic weight, that is, $v=\phi_A+u$ with $u\in\PSH(B,\theta)$ and $u\leq0$, and suppose that $v\leq c\log|x|^2+O(1)$ near $Z$. Write $v_0$ for its expression in the chart $Z_0\neq0$. The obstacle gives $v_0\leq F_0$ on $\{s\leq s_0\}$. On each complex line $z=t\xi$ with $|\xi|=1$, either $v_0(t\xi)\equiv-\infty$, or
$v_0(t\xi)-l\log|t|^2$ is subharmonic on $|t|>\sqrt{s_0}$ and bounded above near infinity. In the coordinate $1/t$ the puncture at infinity is then removable, and the maximum principle bounds this difference by its boundary bound $C$. Varying $\xi$ proves $v\leq F$ on $B$. Consequently $F=\phi_A+u_c$ is the largest such weight, and its contact set $\{F=\phi_A\}$ is $D_c$.

We now return to $X$. By \cref{prop:psh-descent-submersion}, every $\pi^*\theta$-plurisubharmonic function has the form $u\circ\pi$ for a $\theta$-plurisubharmonic function $u$ on $B$. Restriction to the local section $w_1=\cdots=w_r=0$ shows that
$u\circ\pi\preceq\varphi$ implies $u\leq c\log|x|^2+O(1)$ near $Z$.
Conversely, this divisorial bound and the inequality
$|x|^2\leq|x|^2+\sum_j|w_j|^2$ give the domination $u\circ\pi\leq\varphi+O(1)$ near $Y$. A finite cover of $Y$ gives one constant there. On the compact complement of a smaller neighborhood of $Y$, the function $\varphi$ is bounded below and $u\leq0$, so one constant works on all of $X$. The two competitor families therefore correspond to each other, which proves $\mathcal P(\varphi)=u_c\circ\pi$. The pole of this envelope fills $\pi^{-1}(Z)$, although $\varphi$ is finite on $\pi^{-1}(Z)\setminus Y$.

\subsection{Binomial kernels}

For a fixed integer $j\geq0$, set $d=km$, $t_{m,j}=\max\{\ell_m-j,0\}$ and
$\mathcal H_m^{(j)}=H^0(B,\mathcal O(d)(-t_{m,j}Z))$.
The choices $j=0$ and $j=r$ give $W_m$ and $V_m$. Give these spaces the reference norm defined by the Fubini--Study metric, zero continuous weight and the volume $\nu_0=\omega_B^q$. For a binomial random variable $\operatorname{Bin}(d,R)$ with $d$ trials and success probability $R$, their kernels are
\begin{equation}\label{eq:horizontal-binomial-kernel}
 B_{m,j}^0(R)=A_dQ_{m,j}(R),\qquad
 A_d=\binom{d+q}{q},\qquad
 Q_{m,j}(R)=\mathbb P\{\operatorname{Bin}(d,R)\geq t_{m,j}\}.
\end{equation}
To see this, apply the monomial norm formula \eqref{eq:projective-monomial-norm} in degree $d$. Invariance under the coordinate circle actions makes the monomials orthogonal, and the selection keeps exactly those in which $Z_0$ has exponent at least $t_{m,j}$. Summing over the other exponents with the multinomial identity gives \eqref{eq:horizontal-binomial-kernel}. Counting the selected monomials gives
$\dim\mathcal H_m^{(j)}=\binom{d-t_{m,j}+q}{q}=C_cm^q+O(m^{q-1})$.

We next bound the tails from below. Write
$D(v\mid R)=v\log(v/R)+(1-v)\log((1-v)/(1-R))$ for the relative entropy, and put $b=\lfloor da\rfloor=\ell_m$.
If $R\geq b/d$, the binomial mode $\min\{d,\lfloor(d+1)R\rfloor\}$ is at least $b$. Since $t_{m,j}\leq b$, the tail above $t_{m,j}$ contains this mode and is therefore at least $1/(d+1)$. If $0<R<b/d$, the mass at $b$ for the parameter $v=b/d$ is a largest mass, hence at least $1/(d+1)$. For the parameter $R$, the same mass is at least
$(d+1)^{-1}\exp(-dD(v\mid R))$.
Since
\[
 \frac{\partial D(v\mid R)}{\partial v}
       =\log\frac{v(1-R)}{R(1-v)}\geq0\quad(v\geq R)
\]
and $v\leq a$, it is at least $(d+1)^{-1}\exp(-dD(a\mid R))$.
Together with $u_c\leq0$, these cases give
$Q_{m,j}\geq(d+1)^{-1}\exp(mu_c)$, including the cases $b=0$ and $R=1$.
Since $A_d\geq d+1$, the reference kernels satisfy $m^{-1}\log B_{m,j}^0\geq u_c$; at $R=0$ this holds in the extended sense.

For the matching upper asymptotics, fix $0<R<a$. Stirling's formula at the first selected index $t_{m,j}$ gives the logarithmic rate $-kD(a\mid R)$. Since $t_{m,j}/d\to a>R$, the ratios of successive binomial masses above that index are eventually bounded by a number less than one, so the whole tail has the same logarithmic rate. For $R>a$, the binomial variance estimate gives $Q_{m,j}\to1$. At $R=a$, the lower bound by the mode and $Q_{m,j}\leq1$ give the logarithmic limit zero. The functions $m^{-1}\log Q_{m,j}(R)$ increase with $R$, and their limit $u_c$ is continuous on each interval $[\delta,1]$ with $\delta>0$. On a finite partition of such an interval on which the oscillation of $u_c$ is small, monotonicity upgrades pointwise convergence to uniform convergence.

The event that the first $t_{m,j}$ trials succeed gives
\[
 c\log R\leq\frac{t_{m,j}}m\log R
       \leq\frac1m\log Q_{m,j}(R)\leq0.
\]
Since $\log R$ is integrable against every smooth volume, dominated convergence proves global $L^1$ convergence of the functions $m^{-1}\log Q_{m,j}$, and the additional term $m^{-1}\log A_d$ tends to zero. For the measures, the variance estimate gives $Q_{m,j}\to\mathbf1_{D_c}$ away from $R=a$. The pushforward of $\nu_0$ by $R$ is $q(1-R)^{q-1}\,\mathrm{d}R$, by the radial computation in the proof of \cref{thm:hopf-point-potential-limit} applied to $1-R$. In particular the level set $\{R=a\}$ has zero volume. Since $q!A_d/m^q\to k^q$ and these constants are bounded, dominated convergence proves total-variation convergence of $q!m^{-q}B_{m,j}^0\nu_0$ to $\mu_c$, and
\[
 \mu_c(B)=k^q\int_a^1q(1-R)^{q-1}\,\mathrm{d}R=(k-c)^q.
\]

\Needspace{13\baselineskip}
\subsection{Changing the integration volume}

Put $\delta=\min_Bw>0$, and write $B_{m,j}^w$ for the kernel of $\mathcal H_m^{(j)}$ with volume $w\nu_0$. Comparing the two norms gives
\[
 M^{-1}B_{m,j}^0\leq B_{m,j}^w\leq\delta^{-1}B_{m,j}^0.
\]
Together with the preceding subsection, this proves the lower bound \eqref{eq:horizontal-kernel-lower} on the base and the convergence of the functions $m^{-1}\log B_{m,j}^w$. Comparability alone does not identify the measure limit, however.

Let $\Pi_m$ be the orthogonal projection onto $\mathcal H_m^{(j)}$ for the reference $L^2$ norm. In a unit frame at $x$, write $F_{m,x}$ for the evaluation vector of $H^0(B,\mathcal O(d))$ at $x$, and put $K_{m,x}=\Pi_mF_{m,x}$. Formula \eqref{eq:horizontal-binomial-kernel} gives
\[
 \begin{gathered}
 \|F_{m,x}\|_0^2=A_d,\qquad
 \|K_{m,x}\|_0^2=A_dQ_{m,j}(R(x)),\\
 \|F_{m,x}-K_{m,x}\|_0^2=A_d(1-Q_{m,j}(R(x))).
 \end{gathered}
\]
The full evaluation section satisfies
\[
 |F_{m,x}(y)|_{h_A^m}^2=A_d^2\tau(x,y)^d,\qquad
 \tau([Z],[W])=\frac{|\langle Z,W\rangle|^2}{|Z|^2|W|^2}.
\]
Hence $A_d^{-1}|F_{m,x}|^2\nu_0$ is a probability measure concentrating at $x$: outside any neighborhood of $x$ we have $\tau\leq\rho$ for some $\rho<1$, so its mass there is at most $A_d\rho^d\to0$. Continuity of $w$ therefore gives
$\|(w-w(x))F_{m,x}\|_0=o(\sqrt{A_d})$ for each fixed $x$.

On $\mathcal H_m^{(j)}$, let $T_ms=\Pi_m(ws)$. The quadratic form of $T_m$ is the weighted norm, so $\delta I\leq T_m\leq MI$, and
\[
 B_{m,j}^w(x)=\langle T_m^{-1}K_{m,x},K_{m,x}\rangle_0.
\]
For $R(x)>a$, the preceding norm identities give
$\|F_{m,x}-K_{m,x}\|_0=o(\sqrt{A_d})$. Since $(T_m-w(x)I)K_{m,x}=\Pi_m\bigl((w-w(x))K_{m,x}\bigr)$, it follows that
$\|(T_m-w(x)I)K_{m,x}\|_0=o(\sqrt{A_d})$.
Using $\|T_m^{-1}\|\leq\delta^{-1}$, we obtain
\[
 B_{m,j}^w(x)=w(x)^{-1}B_{m,j}^0(x)+o(A_d)
       \qquad(R(x)>a).
\]
For $R(x)<a$, the upper comparison of kernels gives $m^{-q}B_{m,j}^w(x)\to0$.
Finally, $0\leq q!m^{-q}wB_{m,j}^w\leq(M/\delta)q!A_d/m^q$ is uniformly bounded. Dominated convergence away from the null set $\{R=a\}$ proves that
$q!m^{-q}B_{m,j}^w w\nu_0\to\mu_c$ in total variation.

\begin{proof}[Proof of \cref{thm:horizontal-projective-kernels}]
The identification of the selected spaces is \cref{thm:horizontal-center-descent}, and the dimension formula follows from the monomial count. The envelope computation and the kernel calculations for the two endpoint spaces $W_m,V_m$ prove the assertions on the base. Isometric descent gives $B_m^X=B_m^B\circ\pi$.
For $f_c=k^q\mathbf1_{D_c}/w$, fiber integration gives
\[
 \int_X\left|\frac{q!}{m^q}B_m^X-f_c\circ\pi\right|\nu_X
 =\int_B\left|\frac{q!}{m^q}B_m^B-f_c\right|\nu_B.
\]
This proves the total-variation limit on $X$ and the identity $\pi_*\Lambda_{\nu_X}=\mu_c$. The same fiber integration transfers the $L^1$ convergence, and every other smooth positive volume on the compact manifold $X$ is comparable to $\nu_X$. Local uniform convergence pulls back to $X\setminus\pi^{-1}(Z)$. The dimension asymptotics give the probability limits, and $N_m/m^{q+r}=O(m^{-r})$ proves the zero limit in the ambient degree.

For $W_m\subseteq U_m\subseteq V_m$, monotonicity of the kernels gives $B_{W,m}\leq B_{U,m}\leq B_{V,m}$. By \cref{lem:kernel-codimension-comparison}, the positive difference of the $q$-normalized $U_m$ and $W_m$ measures has total mass at most $q!(\dim V_m-\dim W_m)/m^q=O(m^{-1})$. The lower bound and the local uniform convergence for $U_m$ follow from those for the two endpoint spaces by sandwiching. The $L^1$ errors of the endpoint functions bound the error of the intermediate one, which proves global convergence as well.
\end{proof}

\enlargethispage{2\baselineskip}
The limit on the base is independent of $\nu_X$, but its lift is not. If
$\nu_X(\mathrm{d}b,\mathrm{d}x)=\nu_B(\mathrm{d}b)\lambda_b(\mathrm{d}x)$ is the disintegration of $\nu_X$ along the fibers, then
$\Lambda_{\nu_X}(\mathrm{d}b,\mathrm{d}x)=\mu_c(\mathrm{d}b)\lambda_b(\mathrm{d}x)$.
Thus the lifted limit keeps exactly the conditional fiber probabilities over $D_c$.
Conversely, equality of two lifts implies equality of these probabilities for $\mu_c$-almost every $b$: test against $(g\circ\pi)f$, with $g\in C(B)$ and $f$ in a countable dense subset of $C(X)$, and take one common exceptional null set. The reverse implication follows directly by integration.

This dependence occurs in the product example of \cref{ex:horizontal-center-sharp-error}. Keep its selecting function $\varphi=c\log(R_B+R_F)$ and its integer $b\geq c$, and let
$\eta_f=k\omega_B+bf\omega_F$, where $f>0$ is a smooth function on $\mathbb P^r$ with $b\min f>c$. Then
$\eta_f+\ddc\varphi\geq(k-c)\omega_B+(bf-c)\omega_F>0$. The normalized volume of $\eta_f$ and the corresponding limit are
\[
 \nu_{X,f}=\omega_B^q\otimes\frac{f^r\omega_F^r}{\int_{\mathbb P^r}f^r\omega_F^r},
 \qquad
 \Lambda_{\nu_{X,f}}=\mu_c\otimes\frac{f^r\omega_F^r}{\int_{\mathbb P^r}f^r\omega_F^r}.
\]
A nonconstant $f$ changes the lift. For $r\geq2$, it also makes $\eta_f$ nonclosed: $\mathrm{d}\eta_f=b\,\mathrm{d}f\wedge\omega_F$, and wedging with $\omega_F^{r-2}$ reduces the vanishing of this form to the injectivity of $\xi\mapsto\xi\wedge\omega_F^{r-1}$ on one-forms. For $r=1$ this derivative vanishes for degree reasons.

The spaces, the envelope and the kernels do not change when the admissible Hermitian background varies while $\pi,A,h_A,\varphi$ and $\nu_X$ stay fixed. Moving the fiber point of $Y=Z\times\{[1:0:\cdots:0]\}$ to another point of $\mathbb P^r$ does not change them either. For $mc\geq r+1$, however, the constant germ fails to be integrable at the old center and becomes integrable there after the move. Thus the kernels in all degrees together need not determine the multiplier-ideal type of the selecting function.
